\section{Calibration within the smooth physical experiment}
\label{sec:v11-calibration}
The preceding two acquisition theorems supply the gap, area and
multiplier. These quantities cannot be imported from an exact-family
inverse when the unknown experiment has arbitrary smooth remainders.
We instead construct a calibration procedure inside the same selected
physical experiment. Patch labels, a coarse parameter box and common
smoothness bounds remain supplied. The target continues to be the
symmetrized energy invariant, not an arbitrary asymmetric contact graph.

\subsection{Why the integrated flux is stable under calibration}
Let $D<D_+$ be fixed collars. In this subsection assume the relative
$C^3$ estimates on $[0,D_+]$ and a supplied uniform
$C^{3,1}$ bound on $H_{j,b}$ there. These follow on a fixed compact
smooth physical family from the relative theorem with four derivatives.
Suppose a pilot experiment supplies positive values
$\widehat A,\widehat\gamma$ and an upper gap estimate $\widehat g$.
For a programmed window $j\widehat g+d$, put
\[
 \Delta=j(\widehat g-g),\qquad
 R=\frac{\widehat A}{A}
               \frac{\sinh(j\widehat\gamma)}{\sinh(j\gamma)}.
\]
The mean of the scaled bit $2\widehat A\sinh(j\widehat\gamma)Y$
is exactly
\begin{equation}\label{eq:v11-shifted-flux}
                         R H_{j,b}(d+\Delta).
\end{equation}
This is a translation of an integrated flux followed by a constant
rescaling. It has no factor $d^{-2}$. A small calibration error can
therefore be treated as a smooth model error in the Abel discrepancy.

\begin{theorem}[Stability of the plug-in experiment]
\label{thm:v11-plug-in}
Assume $A,\widehat A$ and $\gamma,\widehat\gamma$ lie in fixed
positive compact intervals, with the multiplier interval bounded away
from zero. On a pilot event suppose
\[
 0\le\widehat g-g\le\rho,\qquad
 |\widehat A/A-1|\le\alpha,\qquad
 |\widehat\gamma-\gamma|\le\zeta,
 \qquad j\zeta\le1,\quad j\rho\le D_+-D.
\]
There is a constant $M_H$, determined by the supplied bounds, which
bounds the means in \eqref{eq:v11-shifted-flux} on this event.
Conditionally on any pilot outcome in the event, take fresh independent
bits at the positive grid windows in both orientations. Set
$\widehat c_j=2\widehat A\sinh(j\widehat\gamma)$ and allocate
\begin{equation}\label{eq:v11-plugin-allocation}
 N_{b,k}=\left\lceil
       2\widehat c_j(M_H+t/3)t^{-2}\log\frac{4L}{\eta_1}
                                    \right\rceil.
\end{equation}
The Abel-dictionary estimator, applied to
$Z_{b,k}=\widehat c_j\overline Y_{b,k}$, then satisfies, with
conditional probability at least $1-\eta_1$,
\begin{equation}\label{eq:v11-plugin-error}
 \max_b\|\widehat V_b-\mathcal V_b\|_\infty
 \le C\{h^{m-5/2}+\tau^j+t h^{-5/2}+\delta
                           +\alpha+j\zeta+j\rho\}.
\end{equation}
The allocation has the deterministic bound
$2L+C\widehat c_jh^{-1}t^{-2}\log(4L/\eta_1)$.
Thus calibration errors in this procedure are not multiplied by an
inverse power of the grid spacing. If the pilot event has probability
at least $1-\eta_0$, the unconditional success probability is at
least $1-\eta_0-\eta_1$.
\end{theorem}
\begin{proof}
The true offset of the programmed window is $d+\Delta$. Applying
\eqref{eq:v11-integrated-observation} at this offset proves
\eqref{eq:v11-shifted-flux}. On the fixed multiplier interval,
\[
 \left|\log\frac{\sinh(j\widehat\gamma)}{\sinh(j\gamma)}\right|
 \le j\coth(j\gamma_-)\zeta\le Cj\zeta.
\]
The compact area bounds and $j\zeta\le1$ then give
$|R-1|\le C(\alpha+j\zeta)$ and a fixed upper bound for $R$.
The $C^{3,1}$ bound and the fundamental theorem of calculus through
order two, and the Lipschitz bound at order three, give
\[
 \|H_{j,b}(\cdot+\Delta)-H_{j,b}\|_{C^3([0,D])}
                                                   \le C\Delta.
\]
The upper gap estimate ensures that only nonnegative offsets in the
supplied collar occur, even when the translated function is evaluated
at zero in this analysis. No observation there is made. Together with
the relative law this proves
\begin{equation}\label{eq:v11-calibration-bias}
 \|R H_{j,b}(\cdot+\Delta)-H_b\|_{C^3([0,D])}
                         \le C(\tau^j+\alpha+j\zeta+j\rho).
\end{equation}
It also supplies $M_H$. Conditionally on the pilot, the variance of
a scaled mean is at most $\widehat c_jM_H/N_{b,k}$ and a summand
is bounded by $\widehat c_j$. The same Bernoulli bound and union
argument as before therefore give nodal error at most $t$ with
conditional probability $1-\eta_1$.

Apply \eqref{eq:v11-smooth-bias} to
\eqref{eq:v11-calibration-bias}, and
\eqref{eq:v11-scalar-noise} only to the centered scalar observations.
The dictionary and inverse arguments of
Theorem~\ref{thm:v11-acquisition} give
\eqref{eq:v11-plugin-error}. Summing the ceilings gives the budget;
conditioning and a union bound give the final probability assertion.
\end{proof}

For comparison, normalizing first by $d^2$ would give the mean
$R(1+\Delta/d)^2F_{j,b}(d+\Delta)$, the sensitivity formula in
\cite{A2ReviewV10}. The two calculations concern different processing
of the same bits and are consistent. The present proof avoids the
singular factor before regularization. In the quadratic benchmark
$H(d)=d^2$, translation adds only $2\Delta d+\Delta^2$, which
$\mathscr A$ annihilates exactly. For a general smooth physical flux,
\eqref{eq:v11-calibration-bias} provides the required bound instead
of asserting such an exact cancellation.

\subsection{A charged calibration procedure with smooth remainders}
We now specify sufficient uniform hypotheses for obtaining the pilot.
Let the unknown gap lie in a supplied interval $[g_-,g_+]$ of width
$w_0>0$, with $g_->0$ and $w_0\le d_c$. The area and multiplier
lie in supplied boxes
\[
 0<A_-\le A\le A_+<\infty,\qquad
 0<\gamma_-\le\gamma\le\gamma_+<\infty.
\]
For one selected orientation let $P_r(t)$ denote the success
probability of the $r$-flight event in the physical window $t$.
For $r=1,2$ assume
\begin{equation}\label{eq:v11-calibration-model}
 \begin{split}
 P_r(t)&=0\quad(t\le rg),\\
 P_r(rg+d)&=c_r d^2F_r(d),\qquad
 c_r=\{2A\sinh(r\gamma)\}^{-1},\quad0\le d\le d_c,\\
 F_r(0)&=1,\qquad F_r\ge\beta_c>0,\qquad
                       \|F_r\|_{m-1,1}\le B_c.
 \end{split}
\end{equation}
All constants in these conditions are supplied, not the functions
$F_r$. The exact physical leading coefficient and smooth finite-flight
law give these hypotheses after choosing a common local collar on a
compact smooth family. No equality of the two remainders and no
analytic finite-dimensional parameterization are assumed.

\begin{lemma}[Finite-preparation calibration]
\label{lem:v11-calibration}
Under \eqref{eq:v11-calibration-model}, for sufficiently small
$\xi>0$ and $0<\eta_0<1/2$ there is a measurable procedure using
only selected-event bits at flight numbers one and two which outputs
$\widehat g\in[g_-,g_+]$, $\widehat A\in[A_-,A_+]$ and
$\widehat\gamma\in[\gamma_-,\gamma_+]$. With probability at
least $1-\eta_0$,
\begin{equation}\label{eq:v11-calibration-accuracy}
 0\le\widehat g-g\le C\xi^{1+1/m},\qquad
 |\widehat A/A-1|+|\widehat\gamma-\gamma|\le C\xi.
\end{equation}
Its total number of preparations, including failures, is
deterministically at most
\begin{equation}\label{eq:v11-calibration-cost}
 C\xi^{-(2+2/m)}
                   \log\frac{C\log(C/\xi)}{\eta_0}.
\end{equation}
The constants are uniform over the specified smooth physical class.
\end{lemma}
\begin{proof}
We first obtain a one-sided gap bracket. The bounds imply a supplied
$c_*>0$ such that $P_1(g+d)\ge c_*d^2$ on $[0,d_c]$.
Fix a target bracket width $\rho<w_0$. Starting from $[l,u]$ of
width $w$, sample at $q=l+3w/4$. Use
\[
 N(w)=\left\lceil\frac{16}{c_*w^2}
                         \log\frac{2I}{\eta_0}\right\rceil,
 \qquad I=1+\left\lceil
                     \frac{\log(w_0/\rho)}{\log(4/3)}\right\rceil.
\]
If any success occurs, replace the interval by $[l,q]$; otherwise
replace it by $[l+w/2,u]$. A success certifies $g<q$, because a
success at or below the onset has probability zero. On a correct
incoming bracket a false lower update can occur only if
$g\le l+w/2$. In that case $q-g\ge w/4$ and $q-g\le w\le d_c$,
so its probability is at most
$\exp\{-N(w)c_*w^2/16\}\le\eta_0/(2I)$.
The bracket width decreases by a factor either $3/4$ or $1/2$ at
each step, irrespective of the observations. Stop at width at most
$\rho$. There are at most $I$ steps. A union bound over correct
incoming histories gives a correct final bracket with probability
$1-\eta_0/2$. Its upper endpoint $\widehat g$ then satisfies
$0\le\widehat g-g\le\rho$.

The total of $w^{-2}$ over the steps is at most
$\rho^{-2}\sum_{r\ge0}(3/4)^{2r}$: read the widths backward
from the last pre-stopping width, which is larger than $\rho$.
The ceilings add at most $I$. Consequently this stage costs at most
$C\rho^{-2}\log\{C\log(C/\rho)/\eta_0\}$ preparations.
This bound is deterministic even on the exceptional incorrect-bracket
outcomes.

Choose $s=c_s\xi^{1/m}$ and $\rho=c_\rho\xi s$ with fixed
positive constants small enough that $ms+2\rho\le d_c$ and
$\rho\le s/4$. At fresh preparations, for $r=1,2$ and
$1\le k\le m$, sample the window $r\widehat g+ks$. Let
$\overline Y_{r,k}$ be the corresponding proportion and put
\begin{equation}\label{eq:v11-calibration-extrapolation}
 \widetilde c_r=
   \sum_{k=1}^m(-1)^{k-1}\binom mk
                         \frac{\overline Y_{r,k}}{(ks)^2}.
\end{equation}
The weights reproduce the value at zero of every polynomial of degree
at most $m-1$. Taylor's formula with a Lipschitz highest derivative
therefore bounds the extrapolation error for $c_rF_r$ by $Cs^m$.
If the gap bracket is correct, writing $e=\widehat g-g$ gives the
exact mean at the $k$th normalized node as
\[
 c_r\left(1+\frac{re}{ks}\right)^2F_r(ks+re).
\]
Its difference from $c_rF_r(ks)$ is at most $C\rho/s$;
the finite sum of absolute weights depends only on $m$.

For completeness, at these windows $P_r\le Cs^2$ uniformly in
$k$, since $k\le m$ and $0\le re\le2\rho$. A summand of
$\overline Y_{r,k}/(ks)^2$ is bounded by $s^{-2}$ and its
variance is at most $Cs^{-2}$. Thus
\[
 N_{r,k}=\left\lceil C s^{-2}\xi^{-2}
                                  \log\frac{8m}{\eta_0}\right\rceil
\]
with a sufficiently large supplied constant gives simultaneous
normalized nodal error at most a fixed small multiple of $\xi$
with conditional probability $1-\eta_0/2$. This follows from the
same Bernoulli exponential-moment inequality; its denominator is
bounded by $C s^{-2}(1+\xi)$. Hence, on the joint good event,
\[
 |\widetilde c_r-c_r|
                 \le C(s^m+\rho/s+\xi)\le C\xi.
\]

Project each $\widetilde c_r$ onto the known positive interval of
possible $c_r$ values and call the result $\widehat c_r$. Projection
does not increase its distance from the true $c_r$. Since
$c_1/(2c_2)=\cosh\gamma$, project
$\widehat c_1/(2\widehat c_2)$ onto
$[\cosh\gamma_-,\cosh\gamma_+]$ and apply $\arcosh$ to obtain
$\widehat\gamma$. Finally project
$\{2\widehat c_1\sinh\widehat\gamma\}^{-1}$ onto
$[A_-,A_+]$ to obtain $\widehat A$. All these maps are Lipschitz
on the specified boxes, since $\gamma_->0$. They give
\eqref{eq:v11-calibration-accuracy}. The threshold decisions, finite
sums, and projections are measurable.

The second stage costs $C s^{-2}\xi^{-2}\log(8m/\eta_0)$.
With the selected $s,\rho$, both $s^{-2}\xi^{-2}$ and
$\rho^{-2}$ have order $\xi^{-(2+2/m)}$. Adding the two stages
proves \eqref{eq:v11-calibration-cost}, including all ceilings and
unsuccessful preparations.
\end{proof}

\subsection{The full physical experiment with unknown normalization}
\begin{theorem}[Self-calibrated recovery of boundary profiles]
\label{thm:v11-self-calibrated}
Consider a class of physical channels satisfying the general relative
law, the supplied parameter boxes and calibration conditions
\eqref{eq:v11-calibration-model}, common profile bounds
$\mathcal K_m(B,\beta;D_+)$, and common relative $C^3$ and finite
$C^{3,1}$ flux bounds on $[0,D_+]$. Fix $0<D<D_+$. Only the
selected patch labels, these common bounds and the coarse gap bracket
are known; the actual $g,A,\gamma$ are unknown.

For sufficiently small $\varepsilon>0$ and $0<\eta<1/2$, a
measurable procedure based exclusively on finitely many selected-event
bits recovers both full profiles on $[0,D]$ and predicts the odd law
with uniform error at most $\varepsilon$, with probability at least
$1-\eta$. Its deterministic total preparation count is at most
\begin{equation}\label{eq:v11-self-calibrated-budget}
 C\varepsilon^{-\left(2+\frac6{m-5/2}
                       +\frac{\gamma_+}{|\log\tau|}\right)}
                         \log\frac{C}{\eta\varepsilon}.
\end{equation}
The pilot uses flight numbers one and two; the profile stage uses one
even flight number of order $\log(1/\varepsilon)$. No exact-law
oracle or finite-dimensional remainder model is used. The conclusion
is a uniform sufficient bound, not a minimax lower bound.
\end{theorem}
\begin{proof}
Choose $j,h,t,\delta$ as in Corollary~\ref{cor:v11-budget}, with
constants small enough to accommodate the plug-in error. In the pilot
lemma take $\xi=c\varepsilon/j$ and $\eta_0=\eta/2$.
On its good event, $\alpha\le C\xi$, $\zeta\le C\xi$ and
$\rho\le C\xi^{1+1/m}$. Thus
\[
 \alpha+j\zeta+j\rho\le Cc\varepsilon
\]
for sufficiently small $\varepsilon$. The same choice ensures
$j\zeta\le1$ and $j\rho\le D_+-D$. The latter fixed margin is
why the theorem is stated on a prescribed subcollar. It never loses
the energy origin. Apply Theorem~\ref{thm:v11-plug-in} with fresh
observations, $\eta_1=\eta/2$, and a finite dictionary accurate
enough also for the $2B$ odd-law factor. The union bound and
\eqref{eq:v11-plugin-error} give the error assertion.

The pilot estimates are projected into the supplied boxes even on
bad outcomes. Consequently
$\widehat c_j\le2A_+\sinh(j\gamma_+)$ always, and the profile
allocation has the deterministic bound in
\eqref{eq:v11-self-calibrated-budget}. Before absorption, the pilot
adds at most
\[
 C(j/\varepsilon)^{2+2/m}
               \log\frac{C\log(Cj/\varepsilon)}{\eta}.
\]
Since $j\le C\log(C/\varepsilon)$ and
\[
 2+\frac6{m-5/2}+\frac{\gamma_+}{|\log\tau|}>2+\frac2m,
\]
this additional cost is absorbed into
\eqref{eq:v11-self-calibrated-budget}, uniformly for small
$\varepsilon$ and $0<\eta<1/2$. All observations are independent
fresh preparations conditionally on the chosen windows; their success
or failure is immaterial to the preparation charge.
\end{proof}

The unknown-normalization theorem is a new observation statement, not
an interpretation of the earlier exact-family four-window theorem.
Its smooth-remainder calibration is proved above and its pilot cost is
included. Relative to a calibrated design controlled by the same
supplied parameter box, it does not increase the sufficient accuracy
exponent. This does not replace $\gamma_+$ by an unknown intrinsic
complexity, remove the convergence certificate, or remove patch
labels. Nor does symmetrized energy recovery become recovery of the
two asymmetric branches of a contact graph. Separately observed odd
records test the compatibility predicted by the two recovered even
laws; they are not a third identifying datum for this reconstruction.
